\chapter{2022年新高考II卷}

\section{单选题}

\begin{problem}
已知集合 \(A=\{-1,1,2,4\}\)，\(B=\{x\mid |x-1|\le 1\}\)，则 \(A\cap B=\)（\quad）

\choices
    {\(\{-1,2\}\)}
    {\(\{1,2\}\)}
    {\(\{1,4\}\)}
    {\(\{-1,4\}\)}

\end{problem}

\begin{answer}
B
\end{answer}

\begin{solution}
由
\(|x-1|\le 1\)
得
\(0\le x\le 2\)，
所以
\(B=[0,2]\).
因此
\(A\cap B=\{1,2\}\).
故选 B.
\end{solution}
\begin{problem}
\((2+2\i)(1-2\i)=\)（\quad）

\choices
    {\(-2+4\i\)}
    {\(-2-4\i\)}
    {\(6+2\i\)}
    {\(6-2\i\)}

\end{problem}

\begin{answer}
D
\end{answer}

\begin{solution}
直接计算得
\((2+2\i)(1-2\i)=2-4\i+2\i-4\i^2=6-2\i\).
故选 D.
\end{solution}
\begin{problem}
图 1 是中国古代建筑中的举架结构，\(AA',BB',CC',DD'\) 是桁，相邻桁的水平距离称为步，垂直距离称为举. 图 2 是某古代建筑屋顶截面的示意图，其中 \(DD_1,CC_1,BB_1,AA_1\) 是举，\(OD_1,DC_1,CB_1,BA_1\) 是相等的步，相邻桁的举步之比分别为 \(\frac{DD_1}{OD_1}=0.5\)，\(\frac{CC_1}{DC_1}=k_1\)，\(\frac{BB_1}{CB_1}=k_2\)，\(\frac{AA_1}{BA_1}=k_3\).
已知 \(k_1,k_2,k_3\) 成公差为 \(0.1\) 的等差数列，且直线 \(OA\) 的斜率为 \(0.725\)，则 \(k_3=\)（\quad）

\bitmapfigure[max width=0.82\textwidth,max totalheight=4.8cm]{img/2022/national_paper_2/q3_fig1.png}

\choices
    {0.75}
    {0.8}
    {0.85}
    {0.9}

\end{problem}

\begin{answer}
D
\end{answer}

\begin{solution}
设
\(OD_1=DC_1=CB_1=BA_1=1\).
由题意可得
\(CC_1=k_1\)，\(BB_1=k_2\)，\(AA_1=k_3\).
又因为 \(k_1,k_2,k_3\) 成公差为 \(0.1\) 的等差数列，所以
\(k_1=k_3-0.2\)，\(k_2=k_3-0.1\).

直线 \(OA\) 的斜率为
\(\frac{DD_1+CC_1+BB_1+AA_1}{OD_1+DC_1+CB_1+BA_1}\)
即
\(\frac{0.5+(k_3-0.2)+(k_3-0.1)+k_3}{4}=0.725\).
化简得
\(\frac{3k_3+0.2}{4}=0.725\)，
解得
\(k_3=0.9\).
故选 D.
\end{solution}
\begin{problem}
已知向量 \(\bs a=(3,4)\)，\(\bs b=(1,0)\)，\(\bs c=\bs a+t\bs b\)，若 \(\langle \bs a,\bs c\rangle=\langle \bs b,\bs c\rangle\)，则 \(t=\)（\quad）

\choices
    {\(-6\)}
    {\(-5\)}
    {5}
    {6}

\end{problem}

\begin{answer}
C
\end{answer}

\begin{solution}
由
\(\bs c=\bs a+t\bs b=(3+t,4)\).
又因为 \(\langle \bs a,\bs c\rangle=\langle \bs b,\bs c\rangle\)，所以
\(\frac{\bs a\cdot \bs c}{|\bs a||\bs c|}=\frac{\bs b\cdot \bs c}{|\bs b||\bs c|}\).
于是
\(\frac{3(3+t)+4\cdot 4}{5}=\frac{3+t}{1}\).
即
\(\frac{25+3t}{5}=3+t\).
解得
\(t=5\).
故选 C.
\end{solution}
\begin{problem}
有甲、乙、丙、丁、戊 5 名同学站成一排参加文艺汇演，若甲不站在两端，丙和丁相邻，则不同排列方式共有（\quad）

\choices
    {12 种}
    {24 种}
    {36 种}
    {48 种}

\end{problem}

\begin{answer}
B
\end{answer}

\begin{solution}
把丙、丁捆绑成一个元素，则连同乙、戊共 4 个元素，其中甲不能站在两端.

先排丙丁整体、乙、戊这 3 个元素，有
\(3!\)
种排法.

再将甲插入中间两个位置中的一个，有
\(2\)
种插法.

丙、丁内部还可以交换次序，有
\(2\)
种排法.

所以共有
\(3!\times 2\times 2=24\)
种不同排列方式. 故选 B.
\end{solution}
\begin{problem}
若 \(\sin(\alpha+\beta)+\cos(\alpha+\beta)=2\sqrt2\cos\left(\alpha+\frac{\pi}{4}\right)\sin\beta\)，则（\quad）

\choices
    {\(\tan(\alpha-\beta)=1\)}
    {\(\tan(\alpha+\beta)=1\)}
    {\(\tan(\alpha-\beta)=-1\)}
    {\(\tan(\alpha+\beta)=-1\)}

\end{problem}

\begin{answer}
C
\end{answer}

\begin{solution}
由和角公式得
\(\sin\alpha\cos\beta+\cos\alpha\sin\beta+\cos\alpha\cos\beta-\sin\alpha\sin\beta =2(\cos\alpha-\sin\alpha)\sin\beta\).
整理得
\(\sin\alpha\cos\beta-\cos\alpha\sin\beta+\cos\alpha\cos\beta+\sin\alpha\sin\beta=0\)，
即
\(\sin(\alpha-\beta)+\cos(\alpha-\beta)=0\).
所以
\(\tan(\alpha-\beta)=-1\).
故选 C.
\end{solution}
\begin{problem}
已知正三棱台的高为 1，上、下底面边长分别为 \(3\sqrt3\) 和 \(4\sqrt3\)，其顶点都在同一球面上，则该球的表面积为（\quad）

\choices
    {\(100\pi\)}
    {\(128\pi\)}
    {\(144\pi\)}
    {\(192\pi\)}

\end{problem}

\begin{answer}
A
\end{answer}

\begin{solution}
设正三棱台上、下底面外接圆半径分别为 \(r_1,r_2\)，球半径为 \(R\).

因为正三角形外接圆半径为边长除以 \(\sqrt3\)，所以
\(r_1=\frac{3\sqrt3}{\sqrt3}=3\)，\(r_2=\frac{4\sqrt3}{\sqrt3}=4\).

\bitmapfigure[max width=4.2cm,max totalheight=4.8cm]{img/2022/national_paper_2/q7_solution_fig1.png}

设球心到上、下底面的距离分别为 \(d_1,d_2\). 则
\(d_1=\sqrt{R^2-9}\)，\(d_2=\sqrt{R^2-16}\).
又因为两底面间距离为 1，所以
\(|d_1-d_2|=1\)
或
\(d_1+d_2=1\).
由题意可知符合的是
\(\sqrt{R^2-9}-\sqrt{R^2-16}=1\).
解得
\(R^2=25\).
故球的表面积为
\(S=4\pi R^2=100\pi\).
故选 A.
\end{solution}
\begin{problem}
已知函数 \(f(x)\) 的定义域为 \(\R\)，且 \(f(x+y)+f(x-y)=f(x)f(y)\)，\(f(1)=1\)，则 \(\sum_{k=1}^{22}f(k)=\)（\quad）

\choices
    {\(-3\)}
    {\(-2\)}
    {0}
    {1}

\end{problem}

\begin{answer}
A
\end{answer}

\begin{solution}
令 \(x=1\)，\(y=0\)，得
\(2f(1)=f(1)f(0)\)，
由 \(f(1)=1\) 得
\(f(0)=2\).

令 \(x=0\)，得
\(f(y)+f(-y)=f(0)f(y)=2f(y)\)，
所以
\(f(-y)=f(y)\)，
即 \(f\) 为偶函数.

再令 \(y=1\)，得
\(f(x+1)+f(x-1)=f(x)\).
从而
\(f(x+2)=-f(x-1)=f(x-4)\)，
所以 \(f(x+6)=f(x)\)，即函数周期为 6.

又由递推可得
\(f(2)=f(1)-f(0)=-1\)，\(f(3)=f(2)-f(1)=-2\)，
\(f(4)=f(-2)=f(2)=-1\)，\(f(5)=f(-1)=f(1)=1\)，\(f(6)=f(0)=2\).
于是一个周期内
\(f(1)+f(2)+f(3)+f(4)+f(5)+f(6)=0\).

因为 \(22=6\times 3+4\)，所以
\(\sum_{k=1}^{22}f(k)=f(1)+f(2)+f(3)+f(4)=1-1-2-1=-3\).
故选 A.
\end{solution}
\section{多选题}

\begin{problem}
已知函数 \(f(x)=\sin(2x+\varphi)\)（\(0<\varphi<\pi\)） 的图像关于点 \(\left(\frac{2\pi}{3},0\right)\) 中心对称，则（\quad）

\choices
    {\(f(x)\) 在区间 \(\left(0,\frac{5\pi}{12}\right)\) 单调递减}
    {\(f(x)\) 在区间 \(\left(-\frac{\pi}{12},\frac{11\pi}{12}\right)\) 有两个极值点}
    {直线 \(x=\frac{7\pi}{6}\) 是曲线 \(y=f(x)\) 的对称轴}
    {直线 \(y=\frac{\sqrt3}{2}-x\) 是曲线 \(y=f(x)\) 的切线}

\end{problem}

\begin{answer}
AD
\end{answer}

\begin{solution}
因为图像关于点 \(\left(\frac{2\pi}{3},0\right)\) 中心对称，所以
\(2\cdot \frac{2\pi}{3}+\varphi=k\pi\)（\(k\in\Z\)）.
由 \(0<\varphi<\pi\) 可得
\(\varphi=\frac{2\pi}{3}\).
所以
\(f(x)=\sin\left(2x+\frac{2\pi}{3}\right)\).

对 A，由
\(\frac{\pi}{2}<2x+\frac{2\pi}{3}<\frac{3\pi}{2}\)
解得
\(-\frac{\pi}{12}<x<\frac{5\pi}{12}\).
故在 \(\left(0,\frac{5\pi}{12}\right)\) 上单调递减，A 正确.

对 B，在区间 \(\left(-\frac{\pi}{12},\frac{11\pi}{12}\right)\) 内，函数只经历一个波峰或波谷，故只有一个极值点，B 错误.

对 C，令
\(2x+\frac{2\pi}{3}=k\pi+\frac{\pi}{2}\)，
得
\(x=\frac{k\pi}{2}-\frac{\pi}{12}\)，
所以 \(x=\frac{7\pi}{6}\) 不是其对称轴，C 错误.

对 D，
\(f'(x)=2\cos\left(2x+\frac{2\pi}{3}\right)\).
当 \(x=0\) 时，
\(f(0)=\sin\frac{2\pi}{3}=\frac{\sqrt3}{2}\)，\(f'(0)=2\cos\frac{2\pi}{3}=-1\).
故过点 \(\left(0,\frac{\sqrt3}{2}\right)\) 的切线方程为
\(y-\frac{\sqrt3}{2}=-x\)，
即
\(y=\frac{\sqrt3}{2}-x\).
故 D 正确.

故选 AD.
\end{solution}
\begin{problem}
已知 \(O\) 为坐标原点，过抛物线 \(C:y^2=2px\)（\(p>0\)）焦点 \(F\) 的直线与 \(C\) 交于 \(A,B\) 两点，其中 \(A\) 在第一象限，点 \(M(p,0)\). 若 \(|AF|=|AM|\)，则（\quad）

\choices
    {直线 \(AB\) 的斜率为 \(2\sqrt6\)}
    {\(|OB|=|OF|\)}
    {\(|AB|>4|OF|\)}
    {\(\angle OAM+\angle OBM<180^\circ\)}

\end{problem}

\begin{answer}
ACD
\end{answer}

\begin{solution}
由抛物线 \(y^2=2px\) 的焦点坐标可知
\(F\left(\frac{p}{2},0\right)\).

又因为 \(|AF|=|AM|\)，所以点 \(A\) 在线段 \(FM\) 的垂直平分线上，故
\(x_A=\frac{\frac{p}{2}+p}{2}=\frac{3p}{4}\).
代入抛物线方程得
\(y_A^2=2p\cdot \frac{3p}{4}=\frac{3p^2}{2}\)，
由于 \(A\) 在第一象限，
\(A\left(\frac{3p}{4},\frac{\sqrt6 p}{2}\right)\).

\bitmapfigure[max width=4.0cm,max totalheight=4.8cm]{img/2022/national_paper_2/q10_solution_fig1.png}

设直线 \(AB\) 的斜率为 \(k\)，则
\(k=\frac{\frac{\sqrt6p}{2}-0}{\frac{3p}{4}-\frac{p}{2}}=2\sqrt6\)，
故 A 正确.

由直线 \(AB\) 过焦点 \(F\)，且斜率为 \(2\sqrt6\)，可求另一交点
\(B\left(\frac{p}{3},-\frac{\sqrt6p}{3}\right)\).
于是
\(|OB|=\sqrt{\left(\frac{p}{3}\right)^2+\left(-\frac{\sqrt6p}{3}\right)^2}=\frac{\sqrt7}{3}p\ne \frac{p}{2}=|OF|\)，
故 B 错误.

由抛物线定义，
\(|AB|=|AM|+|BF|=\frac{3p}{4}+\frac{p}{3}+p=\frac{25p}{12}>2p=4|OF|\)，
故 C 正确.

再有 \(\overrightarrow{OA}\cdot \overrightarrow{OB}=\left(\frac{3p}{4},\frac{\sqrt6p}{2}\right)\cdot \left(\frac{p}{3},-\frac{\sqrt6p}{3}\right)=-\frac{3p^2}{4}<0\)，故 \(\angle AOB\) 为钝角.

又 \(\overrightarrow{MA}\cdot \overrightarrow{MB}=\left(-\frac{p}{4},\frac{\sqrt6p}{2}\right)\cdot \left(-\frac{2p}{3},-\frac{\sqrt6p}{3}\right)=-\frac{5p^2}{6}<0\)，故 \(\angle AMB\) 为钝角.

所以
\(\angle OAM+\angle OBM<180^\circ\).
故 D 正确.

故选 ACD.
\end{solution}
\begin{problem}
如图，四边形 \(ABCD\) 为正方形，\(ED\perp\) 平面 \(ABCD\)，\(FB//ED\)，\(AB=ED=2FB\). 记三棱锥 \(E-ACD\)，\(F-ABC\)，\(F-ACE\) 的体积分别为 \(V_1,V_2,V_3\)，则（\quad）

\bitmapfigure[max width=4.1cm,max totalheight=4.8cm]{img/2022/national_paper_2/q11_fig1.png}

\choices
    {\(V_3=2V_2\)}
    {\(V_3=V_1\)}
    {\(V_3=V_1+V_2\)}
    {\(2V_3=3V_1\)}

\end{problem}

\begin{answer}
CD
\end{answer}

\begin{solution}
设
\(AB=ED=2FB=2a\)，
则
\(FB=a\).

于是
\(V_1=\frac13\cdot ED\cdot S_{\triangle ACD} =\frac13\cdot 2a\cdot \frac12(2a)^2 =\frac{4}{3}a^3\)，
\(V_2=\frac13\cdot FB\cdot S_{\triangle ABC} =\frac13\cdot a\cdot \frac12(2a)^2 =\frac{2}{3}a^3\).

\bitmapfigure[max width=4.2cm,max totalheight=4.8cm]{img/2022/national_paper_2/q11_solution_fig1.png}

连接 \(BD\) 交 \(AC\) 于点 \(M\)，连接 \(EM,FM\). 由图可得
\(BD\perp AC\)，\(ED\perp AC\)，
所以
\(AC\perp \text{平面 }BDEF\).

又
\(BM=DM=\frac12 BD=\sqrt2 a\).
过 \(F\) 作 \(FG\perp DE\) 于 \(G\)，则四边形 \(BDGF\) 为矩形，所以
\(FG=BD=2\sqrt2 a\)，\(EG=a\).
从而
\(EM=\sqrt{(2a)^2+(\sqrt2 a)^2}=\sqrt6 a\)，
\(FM=\sqrt{a^2+(\sqrt2 a)^2}=\sqrt3 a\)，
\(EF=\sqrt{a^2+(2\sqrt2 a)^2}=3a\).
故
\(EM^2+FM^2=EF^2\)，
即
\(EM\perp FM\).
于是
\(S_{\triangle EFM}=\frac12 EM\cdot FM=\frac{3\sqrt2}{2}a^2\).

又
\(AC=2\sqrt2 a\)，
所以
\(V_3=\frac13\cdot AC\cdot S_{\triangle EFM} =\frac13\cdot 2\sqrt2 a\cdot \frac{3\sqrt2}{2}a^2 =2a^3\).

因此
\(V_3=V_1+V_2\)，\(2V_3=3V_1\).
故选 CD.
\end{solution}
\begin{problem}
若 \(x,y\) 满足 \(x^2+y^2-xy=1\)，则（\quad）

\choices
    {\(x+y\le 1\)}
    {\(x+y\ge -2\)}
    {\(x^2+y^2\le 2\)}
    {\(x^2+y^2\ge 1\)}

\end{problem}

\begin{answer}
BC
\end{answer}

\begin{solution}
由
\(x^2+y^2-xy=1\)
得
\((x+y)^2=1+3xy\le 1+3\left(\frac{x+y}{2}\right)^2\).
整理可得
\((x+y)^2\le 4\)，
所以
\(-2\le x+y\le 2\).
故 A 错，B 对.

又由
\(xy\le \frac{x^2+y^2}{2}\)
可得
\(1=x^2+y^2-xy\ge x^2+y^2-\frac{x^2+y^2}{2}=\frac{x^2+y^2}{2}\)，
即
\(x^2+y^2\le 2\).
故 C 对.

同时
\(1=x^2+y^2-xy\le x^2+y^2+\frac{x^2+y^2}{2}=\frac32(x^2+y^2)\)，
故
\(x^2+y^2\ge \frac23\).
因此不一定有 \(x^2+y^2\ge 1\)，D 错.

故选 BC.
\end{solution}
\section{填空题}

\begin{problem}
已知随机变量 \(X\) 服从正态分布 \(N(2,\sigma^2)\)，且 \(P(2<X\le 2.5)=0.36\)，则 \(P(X>2.5)=\) \fillinblank{}.

\end{problem}

\begin{answer}
\(0.14\) 或 \(\frac{7}{50}\)
\end{answer}

\begin{solution}
由正态分布关于均值 2 对称可知
\(P(X>2)=\frac12\).
所以
\(P(X>2.5)=P(X>2)-P(2<X\le 2.5)=0.5-0.36=0.14\).
故答案为 \(0.14\) 或 \(\frac{7}{50}\).
\end{solution}
\begin{problem}
曲线 \(y=\ln|x|\) 过坐标原点的两条切线的方程为 \fillinblank{}，\fillinblank{}.

\end{problem}

\begin{answer}
\(y=\frac{1}{\e}x\)，\(y=-\frac{1}{\e}x\)
\end{answer}

\begin{solution}
\bitmapfigure[max width=6.0cm,max totalheight=4.8cm]{img/2022/national_paper_2/q14_solution_fig1.png}

当 \(x>0\) 时，设切点为 \((x_0,\ln x_0)\)，则
\(y'=\frac1x\)，
故切线方程为
\(y-\ln x_0=\frac1{x_0}(x-x_0)\).
又切线过原点，所以
\(-\ln x_0=\frac1{x_0}(0-x_0)=-1\).
解得
\(x_0=\e\).
于是切线方程为
\(y-1=\frac1{\e}(x-\e)\)，
即
\(y=\frac1{\e}x\).

由图像关于 \(y\) 轴对称可知另一条切线方程为
\(y=-\frac1{\e}x\).
故答案为 \(y=\frac1{\e}x\) 和 \(y=-\frac1{\e}x\).
\end{solution}
\begin{problem}
设点 \(A(-2,3)\)，\(B(0,a)\)，若直线 \(AB\) 关于 \(y=a\) 对称的直线与圆 \((x+3)^2+(y+2)^2=1\) 有公共点，则 \(a\) 的取值范围是 \fillinblank{}.

\end{problem}

\begin{answer}
\(\left[\frac13,\frac32\right]\)
\end{answer}

\begin{solution}
点 \(A(-2,3)\) 关于 \(y=a\) 的对称点为
\(A'(-2,2a-3)\).
因为 \(B(0,a)\) 在对称轴 \(y=a\) 上，所以所求对称直线就是直线 \(A'B\).

其方程为
\(y-a=\frac{a-3}{2}(x-0)\)，
即
\((a-3)x+2y-2a=0\).

圆心为 \((-3,-2)\)，半径为 1. 由圆心到直线的距离不大于半径得
\(\frac{| -3(a-3)-4-2a |}{\sqrt{(a-3)^2+4}}\le 1\).
化简得
\((5-5a)^2\le (a-3)^2+4\)，
进一步解得
\(\frac13\le a\le \frac32\).
故答案为
\(\left[\frac13,\frac32\right]\).
\end{solution}
\begin{problem}
已知直线 \(l\) 与椭圆 \(\frac{x^2}{6}+\frac{y^2}{3}=1\) 在第一象限交于 \(A,B\) 两点，\(l\) 与 \(x\) 轴，\(y\) 轴分别交于 \(M,N\) 两点，且 \(|MA|=|NB|\)，\(|MN|=2\sqrt3\)，则 \(l\) 的方程为 \fillinblank{}.

\end{problem}

\begin{answer}
\(x+\sqrt2 y-2\sqrt2=0\)
\end{answer}

\begin{solution}
\bitmapfigure[max width=5.4cm,max totalheight=4.8cm]{img/2022/national_paper_2/q16_solution_fig1.png}

设 \(AB\) 的中点为 \(E\). 由 \(|MA|=|NB|\) 可知
\(ME=NE\)，
所以 \(E\) 也是线段 \(MN\) 的中点.

设直线 \(AB\) 的方程为
\(y=kx+m\)（\(k<0\)，\(m>0\)）.
则
\(M\left(-\frac{m}{k},0\right)\)，\(N(0,m)\)，\(E\left(-\frac{m}{2k},\frac{m}{2}\right)\).

由
\(|MN|=2\sqrt3\)
得
\(\left(\frac{m}{k}\right)^2+m^2=12\).

又由椭圆弦中点性质，可得
\(k_{OE}\cdot k_{AB}=-\frac12\).
而
\(k_{OE}=\frac{m/2}{-m/(2k)}=-k\)，
故
\((-k)\cdot k=-\frac12\)，
即
\(k^2=\frac12\).
由 \(k<0\) 得
\(k=-\frac{\sqrt2}{2}\).

代回 \(|MN|=2\sqrt3\) 得
\(m=2\).
于是
\(y=-\frac{\sqrt2}{2}x+2\)，
即
\(x+\sqrt2 y-2\sqrt2=0\).
故答案为 \(x+\sqrt2 y-2\sqrt2=0\).
\end{solution}
\section{解答题}

\begin{problem}
已知 \(\{a_n\}\) 为等差数列，\(\{b_n\}\) 是公比为 2 的等比数列，且 \(a_2-b_2=a_3-b_3=b_4-a_4\).
\begin{enumerate}
\item 证明：\(a_1=b_1\)；
\item 求集合 \(\{k\mid b_k=a_m+a_1,\ 1\le m\le 500\}\) 中元素个数.
\end{enumerate}

\end{problem}

\begin{solution}
（1）设等差数列 \(\{a_n\}\) 的公差为 \(d\)，等比数列 \(\{b_n\}\) 的首项为 \(b_1\). 则
\(a_2=a_1+d\)，\(a_3=a_1+2d\)，\(a_4=a_1+3d\)，
\(b_2=2b_1\)，\(b_3=4b_1\)，\(b_4=8b_1\).

由
\(a_2-b_2=a_3-b_3\)
得
\(a_1+d-2b_1=a_1+2d-4b_1\)，
所以
\(d=2b_1\).

再由
\(a_2-b_2=b_4-a_4\)
得
\(a_1+d-2b_1=8b_1-(a_1+3d)\).
即
\(a_1+d-2b_1=4d-(a_1+3d)\).
所以
\(a_1=b_1\).

（2）由（1）知
\(d=2b_1=2a_1\).
由 \(b_k=a_m+a_1\) 知
\(b_1\cdot 2^{k-1}=a_1+(m-1)d+a_1\).
所以
\(b_1\cdot 2^{k-1}=b_1+(m-1)\cdot 2b_1+b_1\)，
即
\(2^{k-1}=2m\)，
即
\(m=2^{k-2}\).
又
\(1\le m\le 500\)
得
\(2\le 2^{k-1}\le 1000\).
所以
\(2\le k\le 10\).
故集合中元素个数为
\(10-2+1=9\).
\end{solution}
\begin{problem}
记 \(\triangle ABC\) 的内角 \(A,B,C\) 的对边分别为 \(a,b,c\)，分别以 \(a,b,c\) 为边长的三个正三角形的面积依次为 \(S_1,S_2,S_3\). 已知 \(S_1-S_2+S_3=\frac{\sqrt3}{2}\)，\(\sin B=\frac13\).
\begin{enumerate}
\item 求 \(\triangle ABC\) 的面积；
\item 若 \(\sin A\sin C=\frac{\sqrt2}{3}\)，求 \(b\).
\end{enumerate}

\end{problem}

\begin{solution}
（1）因为边长为 \(x\) 的正三角形面积为 \(\frac{\sqrt3}{4}x^2\)，所以
\(S_1=\frac{\sqrt3}{4}a^2\)，\(S_2=\frac{\sqrt3}{4}b^2\)，\(S_3=\frac{\sqrt3}{4}c^2\).
由
\(S_1-S_2+S_3=\frac{\sqrt3}{2}\)
得
\(\frac{\sqrt3}{4}(a^2-b^2+c^2)=\frac{\sqrt3}{2}\)，
即
\(a^2+c^2-b^2=2\).

由余弦定理
\(b^2=a^2+c^2-2ac\cos B\)
得
\(a^2+c^2-b^2=2ac\cos B=2\).
所以
\(ac\cos B=1\).

又因为
\(\sin B=\frac13\)，
故
\(\cos B=\sqrt{1-\sin^2 B}=\sqrt{1-\frac19}=\frac{2\sqrt2}{3}\).
于是
\(ac=\frac1{\cos B}=\frac{3\sqrt2}{4}\).

所以
\(S_{\triangle ABC}=\frac12 ac\sin B=\frac12\cdot \frac{3\sqrt2}{4}\cdot \frac13=\frac{\sqrt2}{8}\).

（2）由正弦定理
\(\frac{b}{\sin B}=\frac{a}{\sin A}=\frac{c}{\sin C}\)
可得
\(\frac{b^2}{\sin^2 B}=\frac{ac}{\sin A\sin C}\).
代入
\(ac=\frac{3\sqrt2}{4}\)，\(\sin A\sin C=\frac{\sqrt2}{3}\)，\(\sin B=\frac13\)
得
\(\frac{b^2}{(1/3)^2}=\frac{3\sqrt2/4}{\sqrt2/3}=\frac94\).
故
\(b^2=\frac14\)，
所以
\(b=\frac12\).
\end{solution}
\begin{problem}
在某地区进行流行病学调查，随机调查了 100 位某种疾病患者的年龄，得到如下的样本数据的频率分布直方图：

\begin{enumerate}
\item 估计该地区这种疾病患者的平均年龄（同一组中的数据用该组区间的中点值为代表）；
\item 估计该地区一位这种疾病患者的年龄位于区间 \([20,70)\) 的概率；
\item 已知该地区这种疾病的患病率为 \(0.1\%\)，该地区年龄位于区间 \([40,50)\) 的人口占该地区总人口的 \(16\%\). 从该地区中任选一人，若此人的年龄位于区间 \([40,50)\)，求此人患这种疾病的概率（以样本数据中患者的年龄位于各区间的频率作为患者年龄位于该区间的概率，精确到 0.0001）.
\end{enumerate}

\bitmapfigure[max width=5.0cm,max totalheight=4.8cm]{img/2022/national_paper_2/q19_fig1.png}

\end{problem}

\begin{solution}
（1）用各组中点值估算平均年龄，有
\(\bar x=5\times 0.001\times 10+15\times 0.002\times 10+25\times 0.012\times 10\)
\(+35\times 0.017\times 10+45\times 0.023\times 10+55\times 0.020\times 10\)
\(+65\times 0.017\times 10+75\times 0.006\times 10+85\times 0.002\times 10=47.9\).
故估计平均年龄为 47.9 岁.

（2）区间 \([20,70)\) 对应的频率为
\((0.012+0.017+0.023+0.020+0.017)\times 10=0.89\).
故所求概率估计为 0.89.

（3）设事件 \(B\) 表示“此人年龄位于区间 \([40,50)\)”，事件 \(C\) 表示“此人患这种疾病”.

由条件概率公式，
\(P(C\mid B)=\frac{P(BC)}{P(B)}\).
其中
\(P(BC)=0.1\%\times 0.023\times 10\)，\(P(B)=16\%\).
所以
\(P(C\mid B)=\frac{0.1\%\times 0.023\times 10}{16\%}\approx 0.0014\).
故所求概率约为 0.0014.
\end{solution}
\begin{problem}
如图，\(PO\) 是三棱锥 \(P-ABC\) 的高，\(PA=PB\)，\(AB\perp AC\)，\(E\) 是 \(PB\) 的中点.

\begin{enumerate}
\item 证明：\(OE//\) 平面 \(PAC\)；
\item 若 \(\angle ABO=\angle CBO=30^\circ\)，\(PO=3\)，\(PA=5\)，求二面角 \(C-AE-B\) 的正弦值.
\end{enumerate}

\bitmapfigure[max width=4.6cm,max totalheight=4.8cm]{img/2022/national_paper_2/q20_fig1.png}

\end{problem}

\begin{solution}
（1）连接 \(BO\) 并延长交 \(AC\) 于点 \(D\)，连接 \(OA,PD\).

\bitmapfigure[max width=4.8cm,max totalheight=4.8cm]{img/2022/national_paper_2/q20_solution_fig1.png}

因为 \(PO\) 是三棱锥 \(P-ABC\) 的高，所以
\(PO\perp \text{平面 }ABC\).
于是
\(PO\perp AO\)，\(PO\perp BO\).

又因为 \(PA=PB\)，所以在直角三角形 \(\triangle POA\) 与 \(\triangle POB\) 中，
\(PA=PB\)，\(PO=PO\)，
故
\(\triangle POA\cong \triangle POB\).
从而
\(OA=OB\)，\(\angle OAB=\angle OBA\).

又 \(AB\perp AC\)，故
\(\angle OAB+\angle OAD=90^\circ\)，\(\angle OBA+\angle ODA=90^\circ\).
所以
\(\angle OAD=\angle ODA\)，
从而
\(OA=OD\).
因此
\(OA=OB=OD\)，
即 \(O\) 为 \(BD\) 的中点.

又 \(E\) 为 \(PB\) 的中点，所以在三角形 \(PBD\) 中，
\(OE//PD\).
而 \(PD\subset \text{平面 }PAC\)，且 \(OE\not\subset \text{平面 }PAC\)，故
\(OE//\text{平面 }PAC\).

（2）过点 \(A\) 作 \(Az//OP\)，建立空间直角坐标系.

\bitmapfigure[max width=7.8cm,max totalheight=4.8cm]{img/2022/national_paper_2/q20_solution_fig2.png}

由
\(PO=3\)，\(PA=5\)
得
\(OA=\sqrt{PA^2-PO^2}=4\).

又因为
\(\angle OBA=\angle OBC=30^\circ\)，
所以
\(BD=2OA=8\)，\(AD=4\)，\(AB=4\sqrt3\)，\(AC=12\).

可取
\(O(2\sqrt3,2,0)\)，\(B(4\sqrt3,0,0)\)，\(P(2\sqrt3,2,3)\)，\(C(0,12,0)\)，
于是
\(E\left(3\sqrt3,1,\frac32\right)\)，
从而
\(\overrightarrow{AE}=\left(3\sqrt3,1,\frac32\right)\)，\(\overrightarrow{AB}=(4\sqrt3,0,0)\)，\(\overrightarrow{AC}=(0,12,0)\).

设平面 \(AEB\) 的法向量为 \(\bs n=(x,y,z)\)，则
\(\bs n\cdot \overrightarrow{AE}=0\)，\(\bs n\cdot \overrightarrow{AB}=0\).
可取
\(\bs n=(0,-3,2)\).

设平面 \(AEC\) 的法向量为 \(\bs m=(a,b,c)\)，则
\(\bs m\cdot \overrightarrow{AE}=0\)，\(\bs m\cdot \overrightarrow{AC}=0\).
可取
\(\bs m=(3,0,-6)\).

于是
\(\cos\langle \bs n,\bs m\rangle =\frac{\bs n\cdot \bs m}{|\bs n||\bs m|} =-\frac{4\sqrt3}{13}\)，
其中 \(\theta\) 为二面角 \(C-AE-B\) 的大小.

所以
\(\cos\theta=\frac{4\sqrt3}{13}\)，
\(\sin\theta=\sqrt{1-\cos^2\theta} =\sqrt{1-\left(\frac{4\sqrt3}{13}\right)^2} =\frac{11}{13}\).
故二面角 \(C-AE-B\) 的正弦值为
\(\frac{11}{13}\).
\end{solution}
\begin{problem}
已知双曲线 \(C:\frac{x^2}{a^2}-\frac{y^2}{b^2}=1\)（\(a>0\)，\(b>0\)） 的右焦点为 \(F(2,0)\)，渐近线方程为 \(y=\pm \sqrt3 x\).
\begin{enumerate}
\item 求 \(C\) 的方程；
\item 过 \(F\) 的直线与 \(C\) 的两条渐近线分别交于 \(A,B\) 两点，点 \(P(x_1,y_1)\)，\(Q(x_2,y_2)\) 在 \(C\) 上，且 \(x_1>x_2>0\)，\(y_1>0\). 过 \(P\) 且斜率为 \(-\sqrt3\) 的直线与过 \(Q\) 且斜率为 \(\sqrt3\) 的直线交于点 \(M\). 从下面 \(\circled{1}\)\(\circled{2}\)\(\circled{3}\) 中选取两个作为条件，证明另外一个成立：

\(\circled{1}\) \(M\) 在 \(AB\) 上；\(\circled{2}\) \(PQ//AB\)；\(\circled{3}\) \(|MA|=|MB|\).
\end{enumerate}

\end{problem}

\begin{solution}
（1）由右焦点为 \(F(2,0)\) 得
\(c=2\).
又因为渐近线方程为
\(y=\pm \frac{b}{a}x=\pm \sqrt3 x\)，
所以
\(\frac{b}{a}=\sqrt3\)，
即
\(b=\sqrt3\,a\).
由
\(c^2=a^2+b^2\)
得
\(4=a^2+3a^2=4a^2\)，
故
\(a=1\)，\(b=\sqrt3\).
所以双曲线方程为
\(x^2-\frac{y^2}{3}=1\).

（2）由已知得直线 \(PQ\) 的斜率存在且不为零.

若选择由 \(\circled{1}\)\(\circled{2}\) 推 \(\circled{3}\)，或由 \(\circled{2}\)\(\circled{3}\) 推 \(\circled{1}\)，则由 \(\circled{2}\) 可知直线 \(AB\) 的斜率存在且不为零.

若选择由 \(\circled{1}\)\(\circled{3}\) 推 \(\circled{2}\)，则由 \(\circled{3}\) 可知 \(M\) 为线段 \(AB\) 的中点. 若直线 \(AB\) 的斜率不存在，则 \(AB\) 为过点 \(F(2,0)\) 的竖直直线. 此时由双曲线的对称性可知 \(M\) 在 \(x\) 轴上，且
\(M=F\).
再由双曲线关于 \(x\) 轴对称，可得 \(P,Q\) 关于 \(x\) 轴对称，从而
\(x_1=x_2\)，
与 \(x_1>x_2\) 矛盾.

总之，直线 \(AB\) 的斜率存在且不为零.

设直线 \(AB\) 的斜率为 \(k\)，则直线 \(AB\) 的方程为
\(y=k(x-2)\).
设
\(M(x_0,y_0)\).

则条件 \(\circled{1}\)“\(M\) 在 \(AB\) 上”等价于
\(y_0=k(x_0-2)\)，
即
\(ky_0=k^2(x_0-2)\).

两条渐近线的方程可合并为
\(3x^2-y^2=0\).
联立并消去 \(y\)，可得
\((k^2-3)x^2-4k^2x+4k^2=0\).
设 \(A(x_3,y_3)\)，\(B(x_4,y_4)\)，线段中点为 \(N(x_N,y_N)\)，则
\(x_N=\frac{x_3+x_4}{2}=\frac{2k^2}{k^2-3}\)，
\(y_N=k(x_N-2)=\frac{6k}{k^2-3}\).

则条件 \(\circled{3}\)“\(|MA|=|MB|\)”等价于
\((x_0-x_3)^2+(y_0-y_3)^2=(x_0-x_4)^2+(y_0-y_4)^2\).
移项并利用平方差公式整理得
\((x_3-x_4)\bigl[2x_0-(x_3+x_4)\bigr]+(y_3-y_4)\bigl[2y_0-(y_3+y_4)\bigr]=0\).
又因为
\(\frac{y_3-y_4}{x_3-x_4}=k\)，
所以
\(x_0-x_N+k(y_0-y_N)=0\)，
即
\(x_0+ky_0=\frac{8k^2}{k^2-3}\).

由题意知直线 \(PM\) 的斜率为 \(-\sqrt3\)，直线 \(QM\) 的斜率为 \(\sqrt3\)，故
\(y_1-y_0=-\sqrt3(x_1-x_0)\)，\(y_2-y_0=\sqrt3(x_2-x_0)\).
于是
\(y_1-y_2=-\sqrt3(x_1+x_2-2x_0)\).

直线 \(PM\) 的方程为
\(y=-\sqrt3(x-x_0)+y_0\)，
即
\(y=y_0+\sqrt3x_0-\sqrt3x\).
代入双曲线方程 \(3x^2-y^2-3=0\)，得
\(\bigl(y_0+\sqrt3x_0\bigr)\bigl[2\sqrt3x-(y_0+\sqrt3x_0)\bigr]=3\)，
解得点 \(P\) 的横坐标
\(x_1=\frac{1}{2\sqrt3}\left(\frac{3}{y_0+\sqrt3x_0}+y_0+\sqrt3x_0\right)\).
直线 \(QM\) 的方程为
\(y=\sqrt3(x-x_0)+y_0\)，
即
\(y=y_0-\sqrt3x_0+\sqrt3x\).
代入双曲线方程 \(3x^2-y^2-3=0\)，得
\(\bigl(y_0-\sqrt3x_0\bigr)\bigl[-2\sqrt3x-(y_0-\sqrt3x_0)\bigr]=3\)，
解得点 \(Q\) 的横坐标
\(x_2=-\frac{1}{2\sqrt3}\left(\frac{3}{y_0-\sqrt3x_0}+y_0-\sqrt3x_0\right)\).
从而
\(x_1-x_2=\frac{1}{\sqrt3}\left(\frac{3y_0}{y_0^2-3x_0^2}+y_0\right)\)，
\(x_1+x_2-2x_0=-\frac{3x_0}{y_0^2-3x_0^2}-x_0\).
所以直线 \(PQ\) 的斜率
\(m_{PQ}=\frac{y_1-y_2}{x_1-x_2}=\frac{3x_0}{y_0}\).
因此条件 \(\circled{2}\)“\(PQ//AB\)”等价于
\(m_{PQ}=k\)，
即
\(ky_0=3x_0\).

综上所述，条件 \(\circled{1}\) 是 \(ky_0=k^2(x_0-2)\)，条件 \(\circled{2}\) 是 \(ky_0=3x_0\)，条件 \(\circled{3}\) 是 \(x_0+ky_0=\frac{8k^2}{k^2-3}\).

下面说明任取两个都可推出第三个.

若由 \(\circled{1}\)\(\circled{2}\) 出发，则联立
\(ky_0=k^2(x_0-2)\)，\(ky_0=3x_0\)，
解得
\(x_0=\frac{2k^2}{k^2-3}\)，\(ky_0=\frac{6k^2}{k^2-3}\).
所以
\(x_0+ky_0=\frac{8k^2}{k^2-3}\)，
即 \(\circled{3}\) 成立.

若由 \(\circled{1}\)\(\circled{3}\) 出发，则由
\(ky_0=k^2(x_0-2)\)
代入 \(\circled{3}\) 可得
\(x_0+k^2(x_0-2)=\frac{8k^2}{k^2-3}\)，
解得
\(x_0=\frac{2k^2}{k^2-3}\)，\(ky_0=\frac{6k^2}{k^2-3}=3x_0\)，
故 \(\circled{2}\) 成立.

若由 \(\circled{2}\)\(\circled{3}\) 出发，则由
\(ky_0=3x_0\)
代入 \(\circled{3}\)，得
\(4x_0=\frac{8k^2}{k^2-3}\)，
所以
\(x_0=\frac{2k^2}{k^2-3}\)，\(ky_0=\frac{6k^2}{k^2-3}\).
于是
\(ky_0=k^2\left(\frac{2k^2}{k^2-3}-2\right)=k^2(x_0-2)\)，
故 \(\circled{1}\) 成立.

综上，\(\circled{1}\)\(\circled{2}\)\(\circled{3}\) 中任取两个作为条件，都可以推出另外一个成立.
\end{solution}
\begin{problem}
已知函数 \(f(x)=x\e^{ax}-\e^x\).
\begin{enumerate}
\item 当 \(a=1\) 时，讨论 \(f(x)\) 的单调性；
\item 当 \(x>0\) 时，\(f(x)<-1\)，求 \(a\) 的取值范围；
\item 设 \(n\in \N^*\)，证明：\(\frac{1}{\sqrt{1^2+1}}+\frac{1}{\sqrt{2^2+2}}+\cdots+\frac{1}{\sqrt{n^2+n}}>\ln(n+1)\).
\end{enumerate}

\end{problem}

\begin{solution}
（1）当 \(a=1\) 时，
\(f(x)=x\e^x-\e^x=(x-1)\e^x\).
于是
\(f'(x)=x\e^x\).
因为 \(\e^x>0\)，所以
\(f'(x)<0\)（\(x<0\)），\(f'(x)>0\)（\(x>0\)）.
故 \(f(x)\) 的减区间为 \((-\infty,0)\)，增区间为 \((0,+\infty)\).

（2）题设
\(f(x)<-1\)（\(x>0\)）
等价于
\(h(x)=x\e^{ax}-\e^x+1<0\)（\(x>0\)）.
并且
\(h(0)=0\).

求导得
\(h'(x)=(1+ax)\e^{ax}-\e^x\).

若 \(a>\frac12\)，则
\(g(x)=h'(x)\)，\(g'(x)=(2a+a^2x)\e^{ax}-\e^x\).
\(g(0)=0\)，\(g'(0)=2a-1>0\).
因此在 \(0\) 的某个右邻域内有 \(g(x)>g(0)=0\)，即 \(h'(x)>0\). 于是该右邻域内
\(h(x)>h(0)=0\)，
与题意矛盾，所以
\(a\le \frac12\).

下面证明当 \(a\le \frac12\) 时题意成立.

当 \(0<a\le \frac12\) 时，由
\(S(t)=\ln(1+t)-t\)（\(t>0\)），
可得
\(S'(t)=\frac{1}{1+t}-1=-\frac{t}{1+t}<0\).
所以
\(S(t)<S(0)=0\)，
即
\(\ln(1+t)<t\)（\(t>0\)）
从而
\(\ln(1+ax)<ax\).
于是
\((1+ax)\e^{ax}=\e^{ax+\ln(1+ax)}<\e^{2ax}\le \e^x\)（\(x>0\)），
所以
\(h'(x)<0\).

当 \(a\le 0\) 时，对任意 \(x>0\)，有
\(\e^{ax}\le 1\)，\(ax\e^{ax}\le 0\)，
故
\(h'(x)=\e^{ax}+ax\e^{ax}-\e^x<1-1=0\).

综上，当且仅当
\(a\le \frac12\)
时，对一切 \(x>0\) 有 \(h'(x)<0\)，从而
\(h(x)<h(0)=0\)，
即
\(f(x)<-1\).

（3）在（2）中取
\(a=\frac12\)，
则对任意 \(x>0\) 有
\(x\e^{x/2}-\e^x+1<0\).
令
\(t=\e^{x/2}>1\)，
则
\(x=2\ln t\)，\(\e^x=t^2\).
上式化为
\(2t\ln t-t^2+1<0\)，
即
\(2\ln t<t-\frac1t\)（\(t>1\)）.

取
\(t=\sqrt{\frac{n+1}{n}}\)（\(n\in \N^*\)），
则
\(2\ln t=\ln(n+1)-\ln n\)，
并且
\(t-\frac1t =\sqrt{\frac{n+1}{n}}-\sqrt{\frac{n}{n+1}} =\frac{1}{\sqrt{n(n+1)}}\).
因此
\(\ln(n+1)-\ln n<\frac{1}{\sqrt{n^2+n}}\).

对 \(1,2,\dots,n\) 累加得
\(\ln(n+1) =\sum_{k=1}^{n}\bigl(\ln(k+1)-\ln k\bigr) <\sum_{k=1}^{n}\frac{1}{\sqrt{k^2+k}}\).
即
\(\frac{1}{\sqrt{1^2+1}}+\frac{1}{\sqrt{2^2+2}}+\cdots+\frac{1}{\sqrt{n^2+n}}>\ln(n+1)\).
故不等式成立.
\end{solution}
